\documentclass{article}
\usepackage{iclr2027_conference,times}
\usepackage{amsmath,amssymb,amsthm,booktabs,graphicx,multirow}
\usepackage{hyperref,url}
\usepackage{placeins,float}
\hypersetup{hidelinks}
\newcommand{\E}{\mathbb{E}}
\newcommand{\Prb}{\mathbb{P}}
\newcommand{\ind}{\mathbf{1}}
\newtheorem{proposition}{Proposition}
\newtheorem{theorem}{Theorem}
\newcommand{\wape}{\operatorname{WAPE}}
\input{tables/result_macros.tex}
\input{tables/observable_macros.tex}
\title{When More Predictions Cover Less:\\Objective Mismatch under Ratio Risk}
\author{Anonymous authors}
\begin{document}
\maketitle
\lhead{Draft prepared for ICLR 2027}
\begin{abstract}
Selective prediction under a ratio-risk cap can maximize accepted cases or the denominator mass they cover. We prove that these objectives can yield almost maximally different coverage, even with exact conditional information. Reversal requires exposure variation; positive lower and finite upper exposure bounds limit its magnitude. The acceptance unit matters: ordinary binary precision has a common optimum, whereas whole-example abstention for a fixed multi-label classifier permits reversal. For retail weighted absolute percentage error, low-demand cases with near-unit relative error consume little pooled excess budget. Accepted-set accounting reveals this mechanism in mechanically censored M5 data. A one-time external test confirms a learned-policy contrast selected on an earlier holdout: mixed utility gains 5.16 demand-coverage points and loses 3.35 row points. Both prespecified directional bounds exclude zero, while joint operating checks fail. Denominator regularization changes the preferred fitted ranking. Selective predictors should specify coverage utility, report both coverages, and audit accepted-set budget spending.
\end{abstract}

\section{Introduction}
A selective prediction system must specify both its risk constraint and the value of acceptance. Counting accepted cases gives each case equal value. Counting their exposure, opportunity, or demand mass gives cases different weights. When risk is itself a ratio of accepted loss to accepted weight, the choice of coverage utility becomes part of the learning problem.

We establish a sharp mismatch between these two objectives under a common ratio-risk cap. Nonnegative loss and weight variables yield conditional excess and conditional-ratio rankings. A three-context family makes their case-coverage and mass-coverage gaps almost maximal, even with exact conditional information. The construction identifies the needed heterogeneity; a constant conditional denominator, for example, makes the coverage objectives identical. The acceptance unit matters: a fixed multi-label classifier admits a micro-precision reversal when the whole example is accepted or rejected. Ordinary binary precision with an unrestricted gate has a common optimum for both coverage objectives.

Retail WAPE makes the mechanism concrete. A well-predicted high-volume series creates pooled error slack that pays for many poor low-volume predictions. The added rows can have near-unit relative error yet negligible budget cost. CENSORCAST provides matched point forecasts and heads, accepted-set budget accounting, a frozen item-holdout comparison, and a later independent external test of two fixed utility policies. Further analyses examine plug-in score sensitivity, mixed utility, group constraints, and likelihood-based learning from censored sales. These distinguish a population objective mismatch from errors in estimating and transferring a policy.

\section{Related work and scope}
\paragraph{Count-data evaluation and censoring.}
\citet{kolassa2016count} explains why absolute-error measures, including weighted MAPE, can be unsuitable for retail count forecasts. Absolute loss elicits a median; when $\Pr(Y=0)\geq1/2$, zero is a median. Our new object is the acceptance budget: a poor relative-error row may still be inexpensive to accept. M5 established strong tabular forecasting baselines \citep{makridakis2022m5}; we use LightGBM \citep{ke2017lightgbm}. FreshRetailNet provides stockout annotations \citep{wang2025fresh}, and a reproducible forecasting pipeline studies chronological features, censoring, and inventory optimization \citep{azar2026pipeline}. Censored newsvendor work addresses missing demand tails and decisions \citep{hssaine2024censored}. Here the target is acceptance-set WAPE on mechanically hidden recorded sales; inventory decisions and naturally unobserved demand are different estimands.

\paragraph{Ratio metrics and selective prediction.}
Optimization of nondecomposable and ratio-of-linear classification metrics is well established \citep{narasimhan2015complex,bao2020fractional}. Those works optimize prediction performance under the chosen metric. We fix the predictor and compare acceptance utilities under the same ratio constraint, asking how much case coverage can disagree with covered denominator mass.

\paragraph{Selection and risk control.}
Case coverage and its risk trade-off are foundational to selective classification \citep{elyaniv2010foundations,geifman2017selective}. Conditional-loss thresholding is established in reject-option prediction \citep{franc2023reject}, as are group concerns in selective regression \citep{shah2022selective}. Learn then Test and conformal risk control address calibration under explicit assumptions \citep{angelopoulos2021ltt,angelopoulos2024crc}. \citet{joshi2026postprocess} maximize agreement with a baseline decision policy under a chance constraint, ranking contexts by the reduction in conditional violation risk from a fallback action. Our action is forecast acceptance, the constraint is signed expected excess $\E[a(e-r\mu)]\leq0$, and the utility can weight either rows or demand. Accurate high-volume cases can therefore subsidize inaccurate low-volume cases within an accepted set. Conditional thresholding remains the common mathematical tool; our analysis concerns this cross-case budget mechanism and its measurable consequences. Calibration is treated separately.

\section{The acceptance contract}
% Replaces main.tex Section 3.1. Reserve Proposition 1 for Appendix A.
\subsection{Utility mismatch under ratio risk}
Let $X$ be the information available when acceptance is chosen. For
nonnegative integrable loss $L$ and exposure $W$, define
\[
 \ell(X)=\E[L\mid X],\quad w(X)=\E[W\mid X],\quad T=\E W>0,
 \quad a(X)\in[0,1].
\]
Row coverage, exposure coverage, and accepted ratio risk are
\begin{equation}
 c(a)=\E a,\qquad d_W(a)=\frac{\E[aW]}{T},\qquad
 R(a)=\frac{\E[aL]}{\E[aW]}.
 \label{eq:ratio}
\end{equation}
The accepted denominator must be positive. The contract is $R(a)\leq r$
and $c(a)\geq c_0$. Its signed excess is
\begin{equation}
 g_r=L-rW,\qquad m_r=\ell-rw,\qquad
 R(a)\leq r\ \Longleftrightarrow\ \E[a m_r]\leq0.
 \label{eq:excess}
\end{equation}
WAPE is the specialization $L=|Y-f(X)|$, $W=Y$, with $Y,f\geq0$;
then $\ell=e$, $w=\mu$, and $d_W=d$ is demand coverage.

Fixed-row-coverage excess is minimized by accepting the smallest
$m_r$ values, with randomization on ties (Proposition~\ref{prop:threshold}
in Appendix~\ref{app:proofs}). To maximize exposure instead, change measure
to $d\nu=w\,dP_X/T$ after rejecting $w=0,\ell>0$ contexts: then $d_W(a)=\E_\nu a$ and
$R(a)=\E_\nu[a\ell/w]/\E_\nu a$.
\setcounter{proposition}{1}
\begin{proposition}[Exposure-weighted ranking]
\label{prop:weighted}
If a positive-exposure feasible rule exists and the row floor is
nonbinding, exposure coverage is optimized by thresholding $\ell/w$,
with boundary randomization. Reject $w=0,\ell>0$ contexts; $w=\ell=0$
contexts are neutral for this objective. The ranking is independent of
$r$; its threshold is not.
\end{proposition}
Thus the population rankings are $\ell-rw$ for rows and $\ell/w$ for
exposure. A fixed ranking's largest feasible threshold maximizes either
coverage because its acceptance sets are nested; objective comparisons
must therefore also compare ranking families (Appendix~\ref{app:objective}).

For disjoint context sets $K,U$, write
\[
 M_A=\E[\ind_Aw],\qquad B=-\E[\ind_Km_r],\qquad
 M_K>0,\quad B\geq0.
\]
Here $B$ is the excess slack supplied by a common accepted set.
\begin{proposition}[Sharp objective reversal for ratio risk]
\label{prop:lowdemand}
Fix $r>0$, $\rho>r$, and $\varepsilon>0$. If $\ell=\rho w$ on $U$ and
$0<M_U\leq\varepsilon\Pr(U)$, then $R(\ind_U)=\rho$ and
$K\cup U$ satisfies the cap exactly when $(\rho-r)M_U\leq B$.
Accepting $U$ adds $\Pr(U)$ row coverage but at most
$\varepsilon\Pr(U)/T$ exposure coverage.

For every $0\leq c_0<1$ and $\eta>0$, there exists a three-context
law with $w>0$, $0\leq\ell\leq\rho w$, and exact conditional moments
whose unique row- and exposure-optimal policies $a_R,a_W$ satisfy
the cap and the nonbinding row floor, and obey
\[
 c(a_R)-c(a_W)>1-c_0-\eta,\qquad
 d_W(a_W)-d_W(a_R)>1-\eta.
\]
The universal bounds $1-c_0$ and $1$ are jointly sharp. For
$0<r<1$ and $\rho=1$, this sharpness also holds with binary $0\leq L\leq W\leq1$.
\end{proposition}
The sharpness statement ranges over joint laws that can realize the
conditional moments in Appendix~\ref{app:proofs}; it is not a claim
about every fixed metric class. If $W$ is constant, for example,
$d_W=c$ and reversal is impossible. For a fixed multi-label classifier, whole-example abstention yields an exact micro-precision reversal (Appendix~\ref{app:classification}). Ordinary binary precision with an unrestricted gate instead has a common optimum for both coverage objectives. For whole-example micro-precision with $1\leq w\leq m$, each coverage gap is at most $(m-1)/(m+1)$; Appendix~\ref{app:bounded_exposure} sharpens this to a joint bound.

\paragraph{WAPE corollary and estimation.}
For $0<r<1$, setting $\rho=1$, $W=Y$, and $f=0$ on $U$ gives
$e=\mu$: near-zero-demand rows have budget cost $(1-r)\mu$ yet
unit row reward. The three-context law is realizable by deterministic
nonnegative forecasts (Appendix~\ref{app:proofs}), so both sharp bounds
hold for WAPE. A low-exposure share alone cannot bound the forgone
exposure without a positive excess-rate margin; the appendix gives
that bound and the exact budget identity. More generally $|e-\mu|\leq f$.
Exact-head optimality does not imply estimated-ratio robustness.
Section~\ref{sec:objectives} and Appendix~\ref{app:proofs} distinguish
objective alignment from small-denominator error amplification.


\subsection{What censoring does and does not identify}
For natural demand $D$, observe sales $S=\min(D,C)$ and a censoring flag, where $C$ is an availability bound. In a mechanically censored benchmark, the retained pre-censoring target can be used for training and auditing; in natural stockouts it is generally unavailable.

\begin{proposition}[Unbounded censored tails prevent a universal WAPE claim]
\label{prop:identification}
Fix a measurable acceptance probability based on observed information and a bounded nonnegative forecast. Assume the availability bound is nonnegative and integrable and an integrable demand completion exists. Suppose the expected acceptance weight on censored cases is positive and the compatible model class places no upper bound on their demand tail. There exist observationally indistinguishable finite-mean demand completions for which accepted WAPE approaches one. Consequently, observed censored sales alone cannot uniformly establish any contract with $r<1$ over this model class.
\end{proposition}
This statement does not exclude identification under additional structural assumptions. With a valid known bound $D\in[C,U]$ on a censored case, convexity gives the conservative excess envelope
\begin{equation}
 |D-f|-rD \leq
 \max\{|C-f|-rC,\ |U-f|-rU\}.
 \label{eq:envelope}
\end{equation}
Inventory availability $C$ is not itself a demand upper bound $U$. A maximum observed in the development sample is also insufficient to justify a population support bound.

\subsection{Calibration is a separate claim}
\label{sec:calibration}
Independent calibration clusters and known bounded cluster losses permit simultaneous finite-family concentration checks; the supplement states a standard Hoeffding construction (Appendix~\ref{app:calibration}). Our retail counts lack a justified population bound and share store/calendar shocks. Accordingly, development intervals are descriptive; the held-out comparisons use prespecified item-bootstrap approximations, not that theorem.

\section{Method and experimental design}
\subsection{Data, chronology, and the audit boundary}
The M5 split has three disjoint item partitions, with all ten stores of an item kept together. \emph{Development} has 1,829 items (18,290 series) and supplies fitting and threshold selection. The first \emph{item holdout}, called \emph{guardian} in the result files, has 610 items: 42 policies were frozen before its one-time comparison. Its saved predictions subsequently supported exploratory diagnostics. The final \emph{external} partition has 610 other items and is used once for two policies fixed after that exploration (Section~\ref{sec:external}). The first holdout shares development dates and stores; the external primary test uses later forecast origins, still in the same stores. Recorded M5 unit sales are the target $Y$.

After 56 uncensored warmup days, mechanical censoring depends on previous observed sales:
\begin{align}
 L_t&=0.90L_{t-1}+0.10S_{t-1},&
 C_t&=\max\{1,\lceil1.15L_t+0.25\sqrt{L_t+1}\rceil\},&
 S_t&=\min(Y_t,C_t).
 \label{eq:censor}
\end{align}
The previously fixed process hides 30.42\% of recorded demand mass and censors 12.67\% of development rows on days 1314--1913. Initialization is specified in Appendix~\ref{app:repro}.

Horizons are one through seven, with weekly origin $o(t)=t-[1+(t-1)\bmod7]$. Point fits use days 365--1313; selection uses 1314--1433; risk learning uses 1434--1553; two calibration blocks, A and B, use 1554--1673 and 1674--1793; the later development evaluation window (\emph{shadow}) uses 1794--1913. Purging origins that precede fitted-object availability leaves 2,085,060 selective-shadow rows, versus 2,194,800 point-shadow rows. Origin-availability and mask-mutation tests verify the repaired feature constructor (Appendix~\ref{app:repro}).

\subsection{Matched point forecasts and censor-adapter controls}
The reference forecaster fits histogram gradient boosting with Poisson loss to observed sales. Its expectation--maximization (EM) correction repeatedly imputes censored responses using clipped Poisson means conditional on exceeding capacity, then refits; a separate classifier predicts censoring probability. New LightGBM controls share origin-valid features, 1.2 million sampled rows, 300 trees, 63 leaves, learning rate 0.05, and three seeds. Objectives are Poisson or L1 on observed sales, with an uncensored-row-only L1 ablation. For L1 forecast $b$, the adapter reuses the earlier EM prediction $e_{\rm EM}$ and predicted censor probability $\hat p_C$:
\begin{equation}
 f_{\alpha,\gamma}=b+\alpha\hat p_C^\gamma\max(e_{\rm EM}-b,0).
 \label{eq:adapter}
\end{equation}
Only point selection chooses $\alpha,\gamma$. Controls include global/category rescaling and the same ungated correction. The supervised selector comparisons retain the first point seed and $(\alpha,\gamma)=(2,2)$. The sales-and-capacity extension uses the plain observed-sales L1 forecaster without this adapter.

\subsection{Matched selectors and operational checks}
Risk learners share 600,000 sampled risk-training rows, 21 origin-valid features, 220 trees, 31 leaves, and three seeds. Squared-loss heads estimate conditional error $\hat e$, demand $\hat\mu$, or excess. The two-head score is $\hat e-r\hat\mu$ (440 total trees); the matched-budget control uses the first 110 trees of each head. Nonnegative heads are clipped at zero. Forecast-normalized error is $\hat e/\max(f,0.25)$; demand-normalized error is $\hat e/\max(\hat\mu,0.25)$. These names are used throughout.

We require WAPE at most $r=0.6401298$ and row coverage at least 0.35 in both calibration blocks. Each score accepts a ranked prefix. We search all distinct thresholds and choose the feasible threshold with the largest minimum A/B row coverage. Category-specific thresholds follow the same rule; an infeasible category is fully rejected.

The cap is a 15\% reduction from the reference Poisson model's WAPE, 0.7530939208, on all 2,194,800 development point-selection rows (days 1314--1433, before origin purging). The 0.35 row floor retains a five-point buffer above the original 0.30 minimum; it is a study setting, not a retail standard.

The original guardian comparison applies only frozen policies. Subsequent analyses reuse its saved predictions and are explicitly \emph{post-hoc}. They select thresholds and ranking families on development calibration A/B, then apply them to cached guardian outcomes. The exploratory analyses cover objective choice, floors including zero, four head-size pairs, matched Poisson/Tweedie demand heads, and a five-point mixed-utility path. Exact threshold searches are recorded, and saved predictions permit full frontier reconstruction; stored plotting curves are subsampled. Item bootstrap resamples an item's stores and days together; intervals do not include model-fitting uncertainty or remove shared shocks.

\section{Results}
\subsection{More accepted rows can mean less covered demand}
\label{sec:mass}
For a common accepted set $K$ and disjoint added set $U$, let $M_B$ and $E_B$ be the demand and absolute-error mass of set $B$. Excess is additive: $(E_{K\cup U}-rM_{K\cup U})=(E_K-rM_K)+(E_U-rM_U)$. When both masses are positive, the combined set meets the cap exactly when
\begin{equation}
 M_U(R_U-r)\leq M_K(r-R_K).
 \label{eq:budget}
\end{equation}
Thus many inaccurate low-volume rows can consume little of the common set's available error budget. We measure each term through an accepted-set decomposition.

\begin{figure}[t]
\centering\includegraphics[width=\linewidth]{figures/row_vs_demand.png}
\caption{Strong forecaster. Left: three-seed means under the same row-coverage threshold selection; row and demand coverage disagree. Right: first-seed direct-excess versus forecast-normalized-error acceptance sets. The extra direct-only set contains more rows but less demand than the relative-only set. Both panels describe the consumed development shadow.}
\label{fig:mass}
\end{figure}
Direct excess gains 6.08 percentage points in mean row coverage over forecast-normalized error, but loses 5.71 points of demand coverage. At the first seed, their common accepted set covers 73.35\% of rows and has WAPE 0.5958, leaving slack under the 0.64013 cap. The direct-only set adds 267,656 rows (12.84\% of all rows), but only 5.69\% of demand; its WAPE is 1.0026. The relative-only set contains 141,617 rows and 11.45\% of demand, with WAPE 0.7915. Mean forecasts in the two differing sets are 0.039 and 1.498 units. In equation~\eqref{eq:budget}, the common set supplies 97,185.97 units of excess slack. The direct-only set consumes 60,141.11 (61.88\%), versus 50,549.56 (52.01\%) for the relative-only set. Both fit within the pooled budget despite their individual WAPE exceeding the cap. This accounts for the row gain through budget use and demand composition.

\subsection{The two-head control changes the architectural conclusion}
\input{tables/twohead_table.tex}
Table~\ref{tab:twohead} reports six feasible score specifications under the original row objective; the seventh, mean absolute error, has no feasible common threshold.

Table~\ref{tab:twohead} does not support an advantage for direct excess learning. Direct, two-head, and matched-budget two-head row coverage are 86.30\%, 86.35\%, and 86.61\%. Direct learning has slightly lower accepted WAPE (0.6245 versus 0.6268 and 0.6269), so these are nearby trade-offs rather than a dominance or equivalence claim. The coverage ordering changes once across the three fits.

The denominator ablation also matters. Replacing $rf$ by $r\hat\mu$ in a composed excess score raises mean row coverage from 81.47\% to 86.35\%. Conversely, normalizing by $\hat\mu$ raises demand coverage to 89.01\% while accepting 74.17\% of rows. The heads can be reused for any $r$; a six-cap sensitivity is reported without additional fitting. The demand-coverage objective is evaluated explicitly in Section~\ref{sec:objectives}.

\subsection{The frozen comparison on previously held-out items}
\label{sec:heldout}
\input{tables/guardian_comparison_table.tex}
The freeze specifies seven scores, two threshold modes, three seeds, and a first-seed direct-versus-two-head primary contrast. The holdout uses the same A/B and later evaluation dates, without recalibration; no holdout outcome selects an original policy. Its shadow contains 695,400 rows on days 1800--1913.

The coverage trade-off persists (Table~\ref{tab:m5guardian}). Across three fixed fits, direct excess increases row coverage by 6.49 percentage points over forecast-normalized error while reducing demand coverage by 4.36 points. Two-head coverage is 87.41\%, compared with 87.37\% for direct learning, with WAPE 0.6297 versus 0.6280. At the prespecified first seed, direct minus two-head row coverage is $-0.215$ points (paired 95\% item-bootstrap interval $[-0.413,-0.018]$ points); direct WAPE is lower by 0.00141. These are different trade-offs, not an equivalence result or a direct-learning advantage.

\begin{figure}[t]
\centering\includegraphics[width=\linewidth]{figures/guardian_groups.png}
\caption{Previously held-out items, first prespecified seed and frozen pooled thresholds. Dots are shadow WAPE; whiskers extend to the one-sided bootstrap upper bound at nominal tail probability $0.05/48$. Hobbies and Household fail the cap even at their point estimates. All displayed adjusted upper bounds exceed the cap.}
\label{fig:guardiangroups}
\end{figure}

Operating constraints do not transfer jointly. First-seed direct and two-head shadow WAPE are 0.6282 and 0.6296 pooled, but approximately 0.879 for Hobbies and 0.682--0.683 for Household. In the holdout window matching calibration B, even pooled WAPE is 0.6495 and 0.6504, above the 0.64013 cap. All frozen category-specific policies abstain entirely on Hobbies because their development threshold is absent. The prespecified 48-endpoint family covers two primary methods, three periods, four strata, and two bounds, using 10,000 item-bootstrap draws. Every coverage lower bound exceeds 0.35, but every adjusted WAPE upper bound exceeds the cap; shadow pooled upper bounds are 0.6959 and 0.6971. This fails the specified joint operating check. Appendix~\ref{app:heldout} reports all bounds and the access record.

\subsection{Error information limits Hobbies acceptance}
\label{sec:groups}
Replacing only the error head with realized absolute error opens 76.78\% development-shadow Hobbies row coverage; replacing only demand opens no feasible prefix (Appendix Table~\ref{tab:partialoracles}). Both replacements permit 89.15\%. This exposes a samplewise selection-information gap conditional on the fixed forecast. Realized outcomes contain irreducible noise, so these numbers do not establish attainable gains from a perfect conditional learner.

We test feature addition without changing the training sample: adding eight origin-safe intermittency summaries to the shared 600,000-row error head improves Hobbies error MSE from 2.3554 to 2.3370 and positive-day MSE from 7.8403 to 7.7650. No common feasible A/B threshold emerges. Two Hobbies-only heads, with 21/29 features on the same 107,384 rows, instead have MSE 2.4232/2.4399 (Appendix Table~\ref{tab:hobbies_intervention}). The shared-head comparison isolates feature addition; specialization also changes the sample. These tests do not exclude point-forecast misspecification or irreducible error.

A separate group-budget LP uses randomized bins of predicted error and demand, with each category constrained in both A/B blocks. Hobbies has only positive-excess occupied bins (minimum WAPE 0.7290/0.7504), so no nonempty mixture is feasible even without the row floor. A dual certificate verifies this limitation of the fixed score-bin class (Appendix~\ref{app:grouplp}).

Hobbies has 71.48\% zero-demand rows, yet errors on those rows account for only 13.60\% of its absolute error. Positive-demand days alone contribute WAPE 0.7411, above the cap. Intermittency alone therefore does not explain the group failure. Capacity controls also implicate error estimation: enlarging only the error head reduces row coverage, whereas enlarging only demand increases it and WAPE (Appendix~\ref{app:objective}).

\subsection{Objective selection depends on estimated rankings}
\label{sec:objectives}
\input{tables/objective_table.tex}
All 27 score/seed pairs agree between row and demand threshold objectives: 24 yield the same finite threshold and three are infeasible. Choosing among the seven original families on development A/B instead selects direct excess for rows and demand-normalized error for demand. On cached guardian shadow, the latter gains 4.94 demand-coverage points while losing 13.74 row points (Table~\ref{tab:objectives}).

The family comparison in Table~\ref{tab:objectives} uses $\hat e/\max(\hat\mu,0.25)$, not exact $e/\mu$. At floor 0.01, forecast normalization has higher first-seed guardian demand coverage (87.75\% versus 87.53\%); the same ordering holds at 0.05 and 0.1, and reverses at 0.25. With no positive floor, the new calculation gives 87.77\% versus 87.53\%, confirming the limitation. Positive-error/zero-denominator rows rank last, with $0/0$ assigned zero. Population optimality of exact conditional ratios implies no universal ordering of estimated or regularized scores. For $\mu>0$, holding $e$ fixed,
\[
 \frac{\partial(e/\mu)}{\partial\mu}=-\frac{e}{\mu^2},\qquad
 \frac{\partial(e-r\mu)}{\partial\mu}=-r.
\]
More exactly, estimation errors $(\Delta_e,\Delta_\mu)$ change the ratio by $(\mu\Delta_e-e\Delta_\mu)/[\mu(\mu+\Delta_\mu)]$ when $\mu+\Delta_\mu>0$. Small denominators can amplify errors in either head; their covariance also matters. A positive floor limits this sensitivity, while a positive row weight in a mixed utility keeps its denominator away from zero (Appendices~\ref{app:proofs} and~\ref{app:mixed}). This explains a possible instability, not the measured cause of any particular ranking reversal. The floor is reported, not chosen from guardian performance.


Poisson and Tweedie demand heads raise guardian demand coverage to 89.40\% and 89.51\%, but period-B WAPE remains 0.6470 and 0.6467. Head loss and regularization affect empirical results without repairing the joint contract.

\subsection{A constructive path between coverage objectives}
Let $U_\lambda=\lambda c+(1-\lambda)d$, $T=\E Y$, and $w_\lambda=\lambda+(1-\lambda)\mu/T$. For a nonbinding row floor, the same weighted-measure argument orders
\begin{equation}
 s_\lambda=\frac{e-r\mu}{\lambda+(1-\lambda)\mu/T},\qquad 0\leq\lambda\leq1.
 \label{eq:mixed}
\end{equation}
It recovers excess at one endpoint and demand-ratio ordering at the other. We evaluate the fixed grid $\{0,0.25,0.5,0.75,1\}$ with the same first-seed heads, no additional max-denominator regularization, and a common development-only estimate of $T$. All thresholds are chosen on development before application to consumed guardian predictions (Appendix~\ref{app:mixed}).

Appendix Table~\ref{tab:mixed-development} reports the full development A/B path. On the item holdout, at $\lambda=0.25$, shadow row/demand coverage is 84.31\%/88.67\% with WAPE 0.6272; at the composed-excess endpoint $\lambda=1$, it is 87.46\%/83.78\% with WAPE 0.6296. This offers an explicit utility choice. The transferred points need not form an empirical Pareto frontier, and all five still fail period-B pooled WAPE. No mixture is promoted to a confirmed operating policy.

\FloatBarrier
\begin{table}[H]
\centering\small
\caption{One-time external comparison: 610 items, days 1919--1941. Two policies were fixed before opening; thresholds and heads are unchanged. Point WAPE passes the pooled cap, but the 48-endpoint operating check fails.}
\label{tab:external}
\begin{tabular}{lrrr}
\toprule
Fixed policy & Rows (\%) & Demand (\%) & WAPE\\
\midrule
Row utility, $\lambda=1$ & 85.92 & 83.19 & 0.6065 \\
Mixed utility, $\lambda=0.25$ & 82.57 & 88.35 & 0.6081 \\
\bottomrule
\end{tabular}
\end{table}

\subsection{A one-time external confirmation of the fixed contrast}
\label{sec:external}
After inspecting the first-holdout path, we replaced an earlier pending policy pair with $\lambda=1$ and $\lambda=0.25$, then froze their thresholds and first-seed heads before external access (Appendix~\ref{app:external}). The 610 unseen items contribute 140,300 rows on days 1919--1941; all forecast origins follow the last consumed development day. Mixed utility gains 5.16 demand-coverage points and loses 3.35 row points (Table~\ref{tab:external}). Paired one-sided 97.5\% item-bootstrap bounds are $+3.17$ points for the demand difference's lower bound and $-2.48$ for the row difference's upper bound. Both prespecified directions pass. The test concerns this selected learned-policy contrast; it does not test the population theorem, pure-demand optimality, or the full utility path.
All 24 secondary row lower bounds pass, but all 24 adjusted WAPE upper bounds fail. A descriptive decomposition of the same saved masks links the contrast to equation~\eqref{eq:budget}: each exclusive set spends about 30\% of the common slack, while the row-only set has WAPE 1.0046 and carries much less demand (Appendix~\ref{app:external}). This adds mechanism accounting, not a confirmatory endpoint.

\subsection{Observed-sales likelihood does not yet deliver robust control}
\label{sec:observable}
Can the selector avoid pre-censoring training labels? We fit Poisson and negative-binomial (NB) count distributions using observed sales $S$ and availability $C$, holding the plain observed-sales L1 forecast fixed. A capacity hit reveals $Y\geq C$ and contributes a survival probability to the likelihood. The fitted distribution implies both absolute error and demand; A/B thresholds use their ratio, without hidden outcomes. Validation likelihood selects NB shape $\kappa$ (Appendix~\ref{app:observable}).

The selected $\kappa=2$ control passes pooled point estimates but fails Hobbies. Across three fits, shadow row/demand coverage is 48.17\%/54.80\% and realized WAPE is 0.5461, versus implied WAPE 0.6365. Implied error mass exceeds realized error by 20.59\%, while demand mass is 3.46\% higher (Table~\ref{tab:observablemoments}). The aggregate error conservatism explains this ratio gap; the pass does not validate conditional moments.

The first-seed grid $\kappa\in\{1,1.5,2,3,5\}$ exposes instability (Table~\ref{tab:nb-fine-grid}). Shapes 1 and 1.5 admit no threshold; 2 accepts 48.74\% of rows; 3 accepts 79.48\% but violates realized A/B caps; and 5 accepts 99.41\% with shadow WAPE 0.6748. Validation mean NLL ranges from 0.7314 to 0.7347 over shapes 1.5--5. Among four adjusted paired item-bootstrap comparisons, only the NLL-difference interval for $\kappa=3$ relative to selected $\kappa=2$ contains zero. Yet $\kappa=3$ violates both realized A/B caps, which $\kappa=2$ passes at pooled point estimates (Table~\ref{tab:nb-fine-grid}). Thus this likelihood comparison does not resolve a passing shape from a failing one on these calibration checks. The intervals reuse shape-selection data and condition on fitted models.

For an identical-forecast comparison, a supervised 220-tree excess head is truncated to the selected NB's exact row count in each block, with outcome-independent tie randomization. Its shadow WAPE is 0.5321 versus NB's 0.5484, but demand coverage is 52.44\% versus 55.81\%. Neither dominates. This is a descriptive count-matched ranking comparison; Appendix~\ref{app:observable} reports fractional-tie checks. Overall, the likelihood experiment supplies no demonstrated robust operating rule from sales and capacity alone.

\FloatBarrier

\section{Limitations}
The full operating contract fails across periods and categories. The external test concerns two learned policies selected after first-holdout exploration; the remaining utility path is exploratory. Item-bootstrap inference conditions on those fits and does not remove shared calendar/store shocks. Hobbies interventions leave point-forecast misspecification and irreducible outcome noise unresolved. Natural stockouts require justified tail and observation assumptions or independent lost-demand measurements.

Foundation-forecaster comparisons, official M5 metrics, and live inventory trials remain untested. Operational use also requires justified caps, group constraints, and fallback evaluation; covered demand is not profit or service level.

\section{Conclusion}
Under ratio-risk constraints, maximizing accepted cases can sacrifice covered exposure even with exact conditional information. We identify necessary conditions for this mismatch and bound its magnitude given positive lower and finite upper exposure bounds. The retail comparisons connect it to measurable budget spending while exposing estimation and transfer failures. Selective predictors should specify their coverage utility, report both coverages, and audit the risk budget across accepted sets and groups.

\label{end:main}

\clearpage
\section*{AI use statement}
OpenAI Codex assisted with literature retrieval, mathematical formulation and proofs, experimental design, implementation and execution, analysis, simulated peer review, figures, and writing. Retail results use the documented data and prediction caches; synthetic observations appear only in labeled experiments. Automated checks cover input hashes, forecast origins, acceptance masks, historical metrics, budget identities, and rational/LP verification of the theory. The external comparison has one recorded opening, immutable fitted objects, and independent recomputation of all registered statistics and bounds. The supplement records verification scope and execution provenance. Responsibility for the submitted content rests with the authors.

\section*{Ethics statement}
Mechanically hidden recorded sales must not be represented as identified customer demand. Pooled acceptance can concentrate errors in particular product groups; we report those failures and do not recommend deployment from these results. No human-subject intervention, customer profiling, or live replenishment action was performed.

\section*{Reproducibility statement}
The supplement supplies experiment scripts, fixed configurations, models, hashes, result records, and origin/mask tests. Appendix~\ref{app:repro} gives preprocessing, dates, and parameters; Appendix~\ref{app:proofs} gives proofs. The anonymous notebook verifies theory and frozen objects and rebuilds the manuscript. External item sufficient statistics reproduce both primary bounds and all 48 secondary endpoints. Preparation/design source regenerates data from the cited M5 release; consumed caches support the exploratory audits.

\bibliography{references}
\bibliographystyle{iclr2027_conference}
\appendix
\renewcommand{\topfraction}{0.9}
\renewcommand{\bottomfraction}{0.9}
\renewcommand{\textfraction}{0.08}
\renewcommand{\floatpagefraction}{0.85}
\setcounter{topnumber}{4}
\setcounter{totalnumber}{5}
\input{appendix_review4.tex}
\end{document}
